\par\addvspace{12pt}\Needspace{9\baselineskip}

\section{积分的应用}

% ITEM: calculus-example-072
\solutionitem{calculus-example-072}{例5.72}
求\(y=x^2\)与\(y=2-x^2\)围成的面积。再求椭圆\(x^2/a^2+y^2/b^2=1\)内部的面积，利用\(x=au\)的换元与第四章伸缩变换的面积倍率相互核对。
另求\(y=|x|\)与\(y=x^2\)围成的总面积，说明积分区间怎样由交点确定。
\solutionanswer
两抛物线交点横坐标为\(\pm1\)，面积\(\int_{-1}^1(2-2x^2)dx=8/3\)。椭圆面积为\(\int_{-a}^a2b\sqrt{1-x^2/a^2}dx=2ab\int_{-1}^1\sqrt{1-u^2}du=\pi ab\)，与单位圆面积乘伸缩倍率\(ab\)一致。最后两曲线在\(-1,0,1\)相交，\(|x|\ge x^2\)对\(|x|\le1\)成立，总面积\(2\int_0^1(x-x^2)dx=1/3\)。


% ITEM: calculus-exercise-50-01
\solutionitem{calculus-exercise-50-01}{习题5.51（1）}
求\(y=x^3\)与\(y=x\)围成的全部有限区域的总面积，说明为什么不能只计算\(\int_{-1}^1(x-x^3)\,dx\)。
\solutionanswer
交点为\(-1,0,1\)，两块面积相同，总面积\(2\int_0^1(x-x^3)dx=1/2\)。直接在\([-1,1]\)积分奇函数得到零，只表示两个有向面积抵消。


% ITEM: calculus-exercise-50-02
\solutionitem{calculus-exercise-50-02}{习题5.51（2）}
求\(x=y^2\)与\(x=2-y\)围成的面积，比较水平与竖直切片的计算。
\solutionanswer
交点的纵坐标为\(-2,1\)，水平切片宽度\(2-y-y^2\)，故面积\(\int_{-2}^1(2-y-y^2)dy=9/2\)。竖直切片须在\(x=1\)处分段，写为\(\int_0^12\sqrt xdx+\int_1^4(2-x+\sqrt x)dx\)，也为\(9/2\)。水平切片更简洁。


% ITEM: calculus-exercise-50-03
\solutionitem{calculus-exercise-50-03}{习题5.51（3）}
对\(0<a<1\)，求\(y=x^2\)、横轴及直线\(x=a\)、\(x=1\)围成区域的面积，并求\(a\)，使它恰为\([0,1]\)上对应总面积的一半。
\solutionanswer
面积\((1-a^3)/3\)，令其等于\(1/6\)，得\(a=2^{-1/3}\)。


% ITEM: calculus-exercise-50-04
\solutionitem{calculus-exercise-50-04}{习题5.51（4）}
求\(y=\ln x\)、横轴及直线\(x=e\)围成的面积。再求\(y=e^x\)、纵轴及直线\(y=e\)围成的面积，比较两区域关于\(y=x\)的对称关系。
\solutionanswer
第一面积\(\int_1^e\ln xdx=1\)。第二以水平切片求\(\int_1^e\ln y\,dy=1\)，或以竖直切片求\(\int_0^1(e-e^x)dx=1\)。交换坐标正是两区域关于\(y=x\)的反射，面积保持。


% ITEM: calculus-example-073
\solutionitem{calculus-example-073}{例5.73}
质点速度为\(v(t)=t^2-2t\)，\(0\le t\le3\)，初始位置为\(s(0)=1\)。求位置函数、总位移和总路程。再设一个容器的流入率为\(2+t\)、流出率为\(t^2\)，求\([0,2]\)内储量的净变化，并判断储量何时最大。
\solutionanswer
\(s(t)=1+\int_0^t(u^2-2u)du=1+t^3/3-t^2\)。总位移\(s(3)-s(0)=0\)。速度在\((0,2)\)负、\((2,3)\)正，两段位移绝对值均为\(4/3\)，总路程\(8/3\)。容器净变化为\(\int_0^2(2+t-t^2)dt=10/3\)。净流率\((2-t)(t+1)\)在\(0\le t<2\)正，\(t>2\)负，故储量在\(t=2\)最大。


% ITEM: calculus-example-074
\solutionitem{calculus-example-074}{例5.74}
一根细杆位于\([0,L]\)，从左端到位置\(x\)的累计质量为\(M(x)=kx^2\)，其中\(k>0\)。求线密度及区间\([L/2,L]\)内的质量，用密度积分核对。累计质量与线密度记录了什么？
\solutionanswer
线密度是\(\rho=M'=2kx\)。后半段质量\(M(L)-M(L/2)=3kL^2/4\)，也等于\(\int_{L/2}^L2kx\,dx\)。累计质量记录从固定起点到当前位置的总量，线密度记录单位长度的局部质量。


% ITEM: calculus-example-075
\solutionitem{calculus-example-075}{例5.75}
一根细杆占据\([0,L]\)，线密度为\(\rho(x)=k(1+x/L)\)，其中\(k,L>0\)。求总质量与质心位置。再说明同时将所有位置的密度乘以同一个正常数时，质心为什么不变。
\solutionanswer
总质量\(m=\int_0^Lk(1+x/L)dx=3kL/2\)，一阶矩\(\int_0^Lx\rho dx=5kL^2/6\)，故质心\(\bar x=5L/9\)。密度同时乘同一正常数时，质量和一阶矩同倍变化，其比值不变。


% ITEM: calculus-exercise-51-01
\solutionitem{calculus-exercise-51-01}{习题5.52（1）}
设\(v(t)=\cos t\)，\(0\le t\le2\pi\)。求位移、路程和平均速度，并与平均速率比较。
\solutionanswer
位移\(\int_0^{2\pi}\cos tdt=0\)，路程\(\int_0^{2\pi}|\cos t|dt=4\)。平均速度为零，平均速率为\(2/\pi\)。


% ITEM: calculus-exercise-51-02
\solutionitem{calculus-exercise-51-02}{习题5.52（2）}
杆的线密度满足\(\rho(a+b-x)=\rho(x)\)。用积分换元证明质心位于中点。
\solutionanswer
记总质量\(m=\int_a^b\rho>0\)，一阶矩\(J=\int_a^bx\rho(x)dx\)。反射换元给出\(J=\int_a^b(a+b-x)\rho(x)dx=(a+b)m-J\)，所以质心\(J/m=(a+b)/2\)。


% ITEM: calculus-exercise-51-03
\solutionitem{calculus-exercise-51-03}{习题5.52（3）}
设\(w\ge0\)、\(\int_a^bw>0\)。证明\(\int_a^b(f-c)^2w\)在\(c=\overline f_w\)时最小，并解释加权平均与平方偏差之间的关系。
\solutionanswer
记\(\mu=\int fw/\int w\)。展开平方并用\(\int(f-\mu)w=0\)，得\(\int(f-c)^2w=\int(f-\mu)^2w+(c-\mu)^2\int w\)。最后一项非负，且仅在\(c=\mu\)为零，所以加权平均是使加权平方偏差最小的常数。


% ITEM: calculus-exercise-51-04
\solutionitem{calculus-exercise-51-04}{习题5.52（4）}
某弹簧的拉力为\(F(x)=kx\)、\(k>0\)。求将其伸长量从\(a\)增加到\(b\)所需的功，其中\(0\le a<b\)。比较从\(0\)到\(a\)与从\(a\)到\(2a\)的功，解释为何相等位移未必对应相等的功。
\solutionanswer
功为\(\int_a^bkx\,dx=k(b^2-a^2)/2\)。从\(0\)到\(a\)为\(ka^2/2\)，从\(a\)到\(2a\)为\(3ka^2/2\)，后者三倍。后半过程每一点的拉力更大，相同位移的局部贡献并不相同。


% ITEM: calculus-exercise-52-01
\solutionitem{calculus-exercise-52-01}{习题5.53（1） *}
本题椭圆为\(x=a\cos t,y=b\sin t\)，\(a,b>0\)，通常取\(0\le t\le2\pi\)正向描绘一周。
对上述椭圆分别取\(t\in[0,4\pi]\)及反向参数，判断式\textup{(5.47)}的结果，解释覆盖次数与方向的作用。

\par\smallskip\noindent 正文所引关系为：
\begin{equation*}
S=\oint x\,dy=\frac12\oint(x\,dy-y\,dx)
 =\frac12\int_\alpha^\beta(\varphi\psi'-\psi\varphi')\,dt.
\end{equation*}
\solutionanswer
椭圆为\(x=a\cos t,y=b\sin t\)。\([0,4\pi]\)正向走两周，公式结果\(2\pi ab\)；反向走一周得\(-\pi ab\)。只有恰好正向一周时才直接等于几何面积，次数与方向都进入有向积分。


% ITEM: calculus-exercise-52-02
\solutionitem{calculus-exercise-52-02}{习题5.53（2） *}
设\(r(t)\)表示一条正向简单闭曲线。证明平移后面积不变；再用\(x=au+cv\)、\(y=bu+dv\)代入面积公式，核对有向面积乘以\(ad-bc\)的规律。
\solutionanswer
平移为\((x+p,y+q)\)时，\(xdy-ydx\)多出\(p\,dy-q\,dx\)，闭曲线上的积分为零。线性变换则有\(xdy-ydx=(ad-bc)(u\,dv-v\,du)\)，所以有向面积乘\(ad-bc\)，几何面积乘其绝对值。


% ITEM: calculus-example-076
\solutionitem{calculus-example-076}{例5.76 *}
求心形线\(r=a(1+\cos\theta)\)、\(0\le\theta\le2\pi\)围成的面积。再用\(r=2a\cos\theta\)、\(-\pi/2\le\theta\le\pi/2\)计算圆的面积，并与直角坐标下的圆方程核对。
\solutionanswer
心形线面积\(\frac{a^2}2\int_0^{2\pi}(1+\cos\theta)^2d\theta=3\pi a^2/2\)。圆的面积\(2a^2\int_{-\pi/2}^{\pi/2}\cos^2\theta d\theta=\pi a^2\)。后者\(r^2=2ar\cos\theta\)，即\((x-a)^2+y^2=a^2\)，圆半径确为\(a\)。


% ITEM: calculus-exercise-53-01
\solutionitem{calculus-exercise-53-01}{习题5.54（1） *}
求曲线\(r=a\cos2\theta\)在\(-\pi/4\le\theta\le\pi/4\)上所围成一瓣的面积，再由对称性求全部四瓣的总面积。
\solutionanswer
一瓣面积\(\frac{a^2}2\int_{-\pi/4}^{\pi/4}\cos^22\theta d\theta=\pi a^2/8\)，四瓣总面积\(\pi a^2/2\)。


% ITEM: calculus-exercise-53-02
\solutionitem{calculus-exercise-53-02}{习题5.54（2） *}
若区域在每条射线上介于\(r_1(\theta)\)与\(r_2(\theta)\)之间，且\(0\le r_1\le r_2\)，推导其面积公式。用同心圆环核对。
\solutionanswer
同一角度小段的外扇形减内扇形，极限得到\(S=\frac12\int(r_2^2-r_1^2)d\theta\)。同心圆环取常值半径并积分\(0\)至\(2\pi\)，得\(\pi(R^2-r^2)\)。


% ITEM: calculus-example-077
\solutionitem{calculus-example-077}{例5.77 *}
求\(y=x\)与\(y=x^2\)在\([0,1]\)上围成区域绕横轴旋转所得几何体的体积。再用截面积法求底面半径为\(R\)、高为\(h\)的正圆锥体积。
\solutionanswer
旋转后的外半径\(x\)、内半径\(x^2\)，故\(V=\pi\int_0^1(x^2-x^4)dx=2\pi/15\)。圆锥在距顶点\(z\)处的截面半径为\(Rz/h\)，所以\(V=\int_0^h\pi R^2z^2/h^2dz=\pi R^2h/3\)。


% ITEM: calculus-example-078
\solutionitem{calculus-example-078}{例5.78 *}
将\(y=x\)与\(y=x^2\)围成的区域绕纵轴旋转，分别用圆环截面与圆柱壳计算体积。再将它绕\(x=-1\)旋转，说明圆柱壳的半径与高度分别怎样变化。
\solutionanswer
水平截面上\(y\le x\le\sqrt y\)，所以圆环法\(V=\pi\int_0^1(y-y^2)dy=\pi/6\)。圆柱壳法为\(2\pi\int_0^1x(x-x^2)dx=\pi/6\)。改绕\(x=-1\)，半径变为\(x+1\)，高度仍为\(x-x^2\)，体积\(2\pi\int_0^1(x+1)(x-x^2)dx=\pi/2\)。


% ITEM: calculus-exercise-54-01
\solutionitem{calculus-exercise-54-01}{习题5.55（1） *}
将\(y=\sqrt{x}\)、横轴及直线\(x=1\)、\(x=4\)围成的区域绕横轴旋转，求体积。再改绕\(y=-1\)旋转，写出圆环的内、外半径。
\solutionanswer
绕横轴体积\(\pi\int_1^4x dx=15\pi/2\)。绕\(y=-1\)时内半径\(1\)，外半径\(1+\sqrt x\)，体积\(\pi\int_1^4(x+2\sqrt x)dx=101\pi/6\)。


% ITEM: calculus-exercise-54-02
\solutionitem{calculus-exercise-54-02}{习题5.55（2） *}
将\(y=x^2\)、\(y=1\)围成的区域绕纵轴旋转，分别用水平截面和圆柱壳计算。为什么使用圆柱壳时，不能把左右两半区域各自旋转所得体积直接相加？
\solutionanswer
水平截面是半径\(\sqrt y\)的圆，\(V=\pi\int_0^1y dy=\pi/2\)。圆柱壳只取右半区域，\(V=2\pi\int_0^1x(1-x^2)dx=\pi/2\)。左右半区域绕纵轴生成同一个立体，相加会重复计算。


% ITEM: calculus-exercise-54-03
\solutionitem{calculus-exercise-54-03}{习题5.55（3） *}
本题椭球为\(x^2/a^2+y^2/b^2+z^2/c^2\le1\)，\(a,b,c>0\)。
补完上文椭球体积的计算，核对结果在\(a=b=c\)时是否回到球体积。
\solutionanswer
椭球\(x^2/a^2+y^2/b^2+z^2/c^2\le1\)在固定\(x\)处的截面面积为\(\pi bc(1-x^2/a^2)\)。积分\(-a\)至\(a\)，得\(4\pi abc/3\)。\(a=b=c=R\)时为球体积\(4\pi R^3/3\)。


% ITEM: calculus-example-079
\solutionitem{calculus-example-079}{例5.79 *}
用参数公式求半径为\(R\)的圆周长。再求曲线\(x=t^2\)、\(y=t^3\)、\(0\le t\le1\)的弧长，比较直接使用参数与先消去参数的计算。
\solutionanswer
圆取\(x=R\cos t,y=R\sin t\)，速率为\(R\)，长度\(2\pi R\)。第二条曲线的速率为\(\sqrt{4t^2+9t^4}=t\sqrt{4+9t^2}\)，故长度\((13\sqrt{13}-8)/27\)。消参为\(y=x^{3/2}\)，用\(\int_0^1\sqrt{1+9x/4}dx\)得到相同结果；参数式省去了消参，而函数式也并不复杂。


% ITEM: calculus-example-080
\solutionitem{calculus-example-080}{例5.80 *}
求心形线\(r=a(1+\cos\theta)\)、\(0\le\theta\le2\pi\)的周长，并与前面的面积计算比较。化简根式时注意\(\sqrt{\cos^2(\theta/2)}=|\cos(\theta/2)|\)。
\solutionanswer
\(r'=-a\sin\theta\)，故弧长被积函数\(\sqrt{r^2+(r')^2}=2a|\cos(\theta/2)|\)。在\(\pi\)处分段，积分为\(8a\)。面积积累\(r^2/2\)，周长积累\(\sqrt{r^2+(r')^2}\)，前者量纲为长度平方，后者为长度。


% ITEM: calculus-exercise-55-01
\solutionitem{calculus-exercise-55-01}{习题5.56（1） *}
求\(x=y^2\)、\(1\le y\le2\)的弧长，比较以\(x\)为变量和以\(y\)为变量的计算。
\solutionanswer
以\(y\)为变量有\(L=\int_1^2\sqrt{1+4y^2}dy=\sqrt{17}-\sqrt5/2+\frac14\ln\frac{4+\sqrt{17}}{2+\sqrt5}\)。以\(x\)为变量则为\(\int_1^4\sqrt{1+1/(4x)}dx\)，再令\(x=y^2\)回到前式。前者更直接。


% ITEM: calculus-exercise-55-02
\solutionitem{calculus-exercise-55-02}{习题5.56（2） *}
摆线的一拱由\(x=a(t-\sin t)\)、\(y=a(1-\cos t)\)、\(0\le t\le2\pi\)表示。求其弧长。
\solutionanswer
速率\(\sqrt{a^2(1-\cos t)^2+a^2\sin^2t}=2a\sin(t/2)\)，在此参数区间非负。积分得一拱弧长\(8a\)。


% ITEM: calculus-exercise-55-03
\solutionitem{calculus-exercise-55-03}{习题5.56（3） *}
对同一段光滑曲线作保持方向的连续可导参数替换，且替换导数为正。用换元公式证明弧长不变，并与参数曲线切线方向的研究联系起来。
\solutionanswer
若\(t=\varphi(u),\varphi'>0\)，新速率为\(\|r'(\varphi(u))\|\varphi'(u)\)，换元积分即原弧长。新切向量是原切向量的正倍数，方向未变；速度数值改变恰与参数步长的改变抵消。


% ITEM: calculus-exercise-55-04
\solutionitem{calculus-exercise-55-04}{习题5.56（4） *}
一根细线沿四分之一圆周\(x=R\cos t\)、\(y=R\sin t\)、\(0\le t\le\pi/2\)放置，单位弧长的质量为\(\rho(t)=k(1+\sin t)\)。求总质量，并说明为什么不能直接计算\(\int_0^{\pi/2}\rho(t)\,dt\)。
\solutionanswer
\(ds=Rdt\)，故质量\(m=kR\int_0^{\pi/2}(1+\sin t)dt=kR(\pi/2+1)\)。给定的是单位弧长密度，必须乘弧长元，直接按\(dt\)积分会漏去位置变化的尺度\(R\)。


% ITEM: calculus-exercise-55-05
\solutionitem{calculus-exercise-55-05}{习题5.56（5） *}
求极坐标曲线\(r=a\theta\)、\(0\le\theta\le1\)的弧长。把结果与径向距离的总增量比较。
\solutionanswer
\(L=a\int_0^1\sqrt{1+\theta^2}d\theta=\frac a2[\sqrt2+\ln(1+\sqrt2)]\)。它大于径向总增量\(a\)，因为除了向外移动，还沿角度方向移动，弧长因子严格大于\(a\)。


% ITEM: calculus-example-081
\solutionitem{calculus-example-081}{例5.81 *}
对椭圆\(x^2/4+y^2=1\)，求长、短轴端点处的密切圆方程。再求抛物线\(y^2=4x\)顶点处的密切圆，比较圆心与焦点的位置。
\solutionanswer
椭圆参数\((2\cos t,\sin t)\)的曲率为\(2/(4\sin^2t+\cos^2t)^{3/2}\)。长轴端点\((\pm2,0)\)的曲率半径\(1/2\)，圆心\((\pm3/2,0)\)，方程\((x\mp3/2)^2+y^2=1/4\)。短轴端点\((0,\pm1)\)半径\(4\)，圆心\((0,\mp3)\)，方程\(x^2+(y\pm3)^2=16\)。抛物线取\((t^2/4,t)\)，顶点曲率\(1/2\)，半径\(2\)，圆心\((2,0)\)，密切圆\((x-2)^2+y^2=4\)。焦点\((1,0)\)在顶点和密切圆圆心的中点。


% ITEM: calculus-exercise-56-01
\solutionitem{calculus-exercise-56-01}{习题5.57（1） *}
将圆\(x=R\cos t\)、\(y=R\sin t\)代入参数曲率公式，核对\(\kappa=1/R\)。再求\(y=x^2\)在原点及\(x=1\)处的曲率，比较它与常值二阶导数\(2\)。
\solutionanswer
圆的行列式分子绝对值为\(R^2\)，速度立方为\(R^3\)，故曲率\(1/R\)。抛物线曲率\(2/(1+4x^2)^{3/2}\)，在\(0,1\)分别为\(2,2/(5\sqrt5)\)。二阶导数虽恒为\(2\)，斜率不同会改变单位弧长上的转向速度，故曲率不恒定。


% ITEM: calculus-exercise-56-02
\solutionitem{calculus-exercise-56-02}{习题5.57（2） *}
用\(x=a\cosh t\)、\(y=b\sinh t\)求双曲线\(x^2/a^2-y^2/b^2=1\)的曲率。求右顶点处密切圆的圆心与半径，并与椭圆长轴端点比较。
\solutionanswer
行列式\(x'y''-y'x''=-ab\)，所以\(\kappa=ab/(a^2\sinh^2t+b^2\cosh^2t)^{3/2}\)。右顶点半径\(b^2/a\)，圆心\((a+b^2/a,0)\)。椭圆长轴右端的半径也为\(b^2/a\)，但圆心在\((a-b^2/a,0)\)，弯曲方向相反。


% ITEM: calculus-exercise-56-03
\solutionitem{calculus-exercise-56-03}{习题5.57（3） *}
对保持方向的二阶连续可导参数替换\(t=\varphi(u)\)、\(\varphi'>0\)，由式\textup{(5.54)}证明曲率不变。二阶求导产生的哪一项在行列式中消失？

\par\smallskip\noindent 正文所引关系为：
\begin{equation*}
\kappa=\frac{|\theta'|}{v}
 =\frac{|x'y''-y'x''|}{(x'^2+y'^2)^{3/2}}.
\end{equation*}
\solutionanswer
新一阶导数为\(r'\varphi'\)，二阶为\(r''(\varphi')^2+r'\varphi''\)。行列式中\(\det(r'\varphi',r'\varphi'')=0\)，故分子乘\((\varphi')^3\)，分母速度立方也乘同一因子，曲率不变。


% ITEM: calculus-exercise-56-04
\solutionitem{calculus-exercise-56-04}{习题5.57（4） *}
将平面图形整体放大\(c>0\)倍，证明曲率变为原来的\(1/c\)，曲率半径变为原来的\(c\)倍。再用圆变椭圆的例子，说明第四章的一般仿射变换为什么没有同一个曲率倍率。
\solutionanswer
整体放大使一、二阶导数均乘\(c\)，曲率分子乘\(c^2\)、分母乘\(c^3\)，故曲率除以\(c\)，半径乘\(c\)。单位圆经不同方向伸缩成椭圆后，长轴端点曲率\(a/b^2\)、短轴端点\(b/a^2\)不同，不能给一般仿射变换指定一个统一倍率。


% ITEM: calculus-exercise-56-05
\solutionitem{calculus-exercise-56-05}{习题5.57（5） *}
椭圆\(x^2/a^2+y^2/b^2=1\)、\(a>b>0\)在哪些点弯曲最急，哪些点最平缓？求两处曲率的比值，并讨论当\(a/b\)增大时怎样变化。
\solutionanswer
参数曲率\(ab/(a^2\sin^2t+b^2\cos^2t)^{3/2}\)。长轴两端最急，曲率\(a/b^2\)；短轴两端最平缓，曲率\(b/a^2\)。最大与最小曲率之比\((a/b)^3\)，随轴比增大而增大。


% ITEM: calculus-example-082
\solutionitem{calculus-example-082}{例5.82 *}
将上半圆\(x=R\cos t\)、\(y=R\sin t\)、\(0\le t\le\pi\)绕横轴旋转，求球面积。再求其中两平行平面\(x=c_1\)、\(x=c_2\)之间的球带面积，观察它与球带轴向高度的关系。
\solutionanswer
参数速率为\(R\)，到横轴距离\(R\sin t\)，所以球面积\(2\pi R^2\int_0^\pi\sin tdt=4\pi R^2\)。球带由\(dx=-R\sin tdt\)得面积\(2\pi R(c_2-c_1)\)，只依赖轴向高度，\(-R\le c_1<c_2\le R\)。


% ITEM: calculus-example-083
\solutionitem{calculus-example-083}{例5.83 *}
将\(y=\sqrt{x}\)、\(1\le x\le4\)绕横轴旋转，求生成的曲面面积。再写出该曲线绕直线\(y=-1\)旋转时的面积积分，说明哪些因子改变、哪些因子保持不变。
\solutionanswer
曲面面积为\(2\pi\int_1^4\sqrt x\sqrt{1+1/(4x)}dx=\pi\int_1^4\sqrt{4x+1}dx=\pi(17\sqrt{17}-5\sqrt5)/6\)。改绕\(y=-1\)时变为\(2\pi\int_1^4(\sqrt x+1)\sqrt{1+1/(4x)}dx\)：旋转半径增加\(1\)，弧长因子不变。


% ITEM: calculus-exercise-57-01
\solutionitem{calculus-exercise-57-01}{习题5.58（1） *}
用积分求底面半径为\(R\)、高为\(h\)的正圆锥侧面积，再求总表面积。比较体积与侧面积中局部贡献的表达式。
\solutionanswer
取母线\(y=Rx/h\)，\(0\le x\le h\)，侧面积\(2\pi\int_0^h(Rx/h)\sqrt{1+R^2/h^2}dx=\pi R\sqrt{R^2+h^2}\)。总表面积再加\(\pi R^2\)。体积的局部贡献是截面积乘轴向厚度，侧面积则是周长乘斜向弧长。


% ITEM: calculus-exercise-57-02
\solutionitem{calculus-exercise-57-02}{习题5.58（2） *}
将\(y=x\)、\(1\le x\le2\)分别绕横轴、纵轴旋转，求曲面面积，用圆台公式核对。
\solutionanswer
绕横轴、纵轴的旋转半径均为\(x\)，\(ds=\sqrt2dx\)，两者面积均为\(2\pi\sqrt2\int_1^2x dx=3\pi\sqrt2\)。圆台上下半径\(1,2\)，母线长\(\sqrt2\)，侧面积\(\pi(1+2)\sqrt2\)，相合。


% ITEM: calculus-exercise-57-03
\solutionitem{calculus-exercise-57-03}{习题5.58（3） *}
设\(R>r>0\)。圆\(x=R+r\cos t\)、\(y=r\sin t\)、\(0\le t\le2\pi\)绕纵轴旋转生成圆环面，求其面积。解释为何此处应选整圆作为母线，而求球面时只选一个半圆。
\solutionanswer
到纵轴的距离为\(R+r\cos t>0\)，\(ds=r dt\)。面积\(2\pi r\int_0^{2\pi}(R+r\cos t)dt=4\pi^2Rr\)。整圆描绘了圆环面的内外两侧，各处恰覆盖一次；球面若用整圆作母线，会让两半圆旋转生成的球面重复覆盖。


% ITEM: calculus-exercise-57-04
\solutionitem{calculus-exercise-57-04}{习题5.58（4） *}
同一条母线改变参数后，利用换元公式证明旋转面积不变。若沿母线走两遍，积分结果如何变化？
\solutionanswer
设换参函数\(t=\phi(s)\)在\([u,v]\)上单调且连续可导，恰将整个参数区间\([\alpha,\beta]\)走过一遍。由链式法则，换参后的速度为\(\|r'(\phi(s))\|\,|\phi'(s)|\)，故旋转面积为
\[
 2\pi\int_u^v \rho(\phi(s))\|r'(\phi(s))\|\,|\phi'(s)|\,ds
 =2\pi\int_\alpha^\beta \rho(t)\|r'(t)\|\,dt.
\]
递增时直接换元即得等式；递减时\(t\)从\(\beta\)走到\(\alpha\)，而\(|\phi'|=-\phi'\)，负号恰好抵消上下限的反转。若沿母线完整走两遍，每一遍的积分都等于原面积，故积分加倍，几何面积不变。
\par
